% !TeX program = pdflatex
% ALIUS Bulletin standalone facsimile source.
% Compile from the ALIUS-bulletin project root for exact visual fidelity.
% If the shared PDF asset is absent, this file falls back to extracted text.
\providecommand{\ALIUSRootPrefix}{}

\ifdefined\ALIUSIssueBuild
  \IfFileExists{\ALIUSRootPrefix Shared-assets/original-pdfs/Issue06/alius-bulletin-issue-06-koroma-what-is-like-to-be-asleep.pdf}{%
    \includepdf[pages=-,fitpaper=true]{\ALIUSRootPrefix Shared-assets/original-pdfs/Issue06/alius-bulletin-issue-06-koroma-what-is-like-to-be-asleep.pdf}%
  }{%
    \typeout{ALIUS warning: missing source PDF Shared-assets/original-pdfs/Issue06/alius-bulletin-issue-06-koroma-what-is-like-to-be-asleep.pdf}%
  }%
  \expandafter\endinput
\fi

\documentclass[a4paper,11pt]{article}
\usepackage[margin=22mm]{geometry}
\usepackage{pdfpages}
\usepackage[hidelinks]{hyperref}
\usepackage{microtype}
\usepackage{parskip}
\title{What is it like to be asleep ?}
\author{An episode with Matthieu Koroma recorded by Daniel Friedman edited by Lena Coutrot, Daniel Friedman and Matthieu Koroma}
\date{ALIUS Bulletin 6}
\begin{document}
\IfFileExists{\ALIUSRootPrefix Shared-assets/original-pdfs/Issue06/alius-bulletin-issue-06-koroma-what-is-like-to-be-asleep.pdf}{%
  \includepdf[pages=-,fitpaper=true]{\ALIUSRootPrefix Shared-assets/original-pdfs/Issue06/alius-bulletin-issue-06-koroma-what-is-like-to-be-asleep.pdf}%
}{%
  \maketitle
  \section*{Metadata}
  \textbf{Piece type:} podcast\par
  \textbf{DOI:} 10.5281/zenodo.7419360\par
  \textbf{Keywords:} Sleep, Dreaming, Lucid Dreaming, Hypgnagogic Dreaming, Memory, Self, Time\par
  \section*{Abstract}
  In this episode, Matthieu Koroma reviews the recent progress on the scientific study of conscious experiences during sleep.\par
  \section*{Extracted Text}
  \begingroup
  \small
  \noindent What is it like to be asleep ? An episode with Mahieu Koroma Recorded by Daniel Friedman and edited by Lena Coutrot and Mahieu Koroma Cite as: Koroma, M. (2022). What is it like to be asleep ? ALIUS Bullen , 6 , hps://doi.org/10.5281/zenodo.7419360 Mahieu Koroma mahieu.koroma@uliege.be University of Liege, Belgium Abstract In this episode, Mahieu Koroma reviews the recent progress on the scienfic study of conscious experiences during sleep. keywords : Sleep, Dreaming, Lucid Dreaming, Hypnagogic Dreaming, Memory, Self, Time The video and audio version of this episode are available on hps://aliusresearch.org/bullen06-koroma-podcast.html Conscious experiences during sleep have long been considered mysterious and typically dicult to assess as they can only be shared aer awakening. Even the memories that we are le with in the morning are oen faint and sometimes bizarre. So a uestion naturally arises : do memories aer awakening truly reect the type of conscious experiences that we have during sleep ? Even more dramatically, could memories be completely made up to such an extent that we actually do not dream at all during sleep? (Malcolm, 1953; Dennett, 1976). is hypothesis has been debated over the past 50 years but the scientic consensus is now suesting that dream reports are actually faithful (Windt, 2013). One of the strongest arguments in favor of this hypothesis comes from people who have a pathological condition called REM behavior disorder. Because people with REM behavior disorder are unable to inhibit their movements while dreaming, researchers could observe their actions and found that they t remarkably well with the dream experiences that were described aer awakening (Oudiette et al., 2009). ALIUS Bulletin n6 (2022) aliusresearch.org/bulletin Matthieu Koroma - What is it like to be asleep ? 2 Could memories be made up to the such an extent that we actually do not dream at all during sleep? So if we can trust dream reports, what do they tell us? Researchers have studied this uestion systematically by waking up participants at diherent moments of the night and asking them to describe the conscious experience that they recall from their dreams. What they found is that dream contents varied greatly d epending on the time of the night (Casagrande et al., 1996; Siclari et al., 2013). While we are falling asleep, we can experience intense conscious activity that is called hypnagogic dreams. ese dreams oen reect what we have been experiencing during the day. is is known as the Tetris ehect, because doing some actions for several hours, such as playing the video game Tetris, causes images of this game to be replayed during sleep (Stickgold, 2000). Even while rmly asleep, our dreams tend to be composed of images and actions related to our daily activities. But as the night goes on, the more complex our dreams become. One aspect of dreams that has drawn particular attention is our experience of the self (Revonsuo, 2005). Our subjective experience can be in the rst person perspective, embracing the point of view of our body. However, we may also experience dreams in the third person, oating slightly above ourselves. We may even be absent from the dream, adopting a disembodied point of view that simply watches the dream unfold. Another focus on the dream experience has been the perception of time. While a dream report can describe actions taking place over several hours or days, we only dream for a few minutes or hours. However, a good correlation was found between the actual time spent asleep and the duration of dream experiences (Dement \& Wolpert, 1958). An explanation for this fact could be that, in our dreams, we are not so surprised by transitions that involve jump in time, a phenomenon that is also found in books or movies and called narrative ellipses. The dreamless sleep experience has been described in parcular detail by great meditators. ALIUS Bulletin n6 (2022) aliusresearch.org/bulletin Matthieu Koroma - What is it like to be asleep ? 3 Finally, we can also wake up knowing that we experienced something in our dream but without being able to say precisely what it was about . While this may reect a failure of our memory, this could also tap into the existence of conscious experiences that are lacking any content. is phenomenon is known as mind-blanking during wakefulness (Ward \& Wegner, 2012) and white dreams during sleep (Fazekas et al., 2019). is dreamless sleep experience has been described in particular detail by great meditators who have trained themselves to remain conscious even aer falling asleep (Windt et al., 2016). Another trainable skill that has sparkled the interest of researchers as a tool to probe the nature of dreams is lucid dreaming. Lucid dreaming is the ability to become aware of dreaming while in a dream (LaBerge \& Rheingold, 1990). A uestion that has been investigated is to what extent the body of the sleeper is involved when we dream. Researchers found that sleepers' respiratory rhythm slows down if lucid dreamers are asked to stop breathing in their dream (Oudiette et al., 2019). By asking diherent uestions to lucid dreamers while recording their body and brain responses, ve laboratories around the world have now succeeded in communicating in real-time with lucid dreamers (Konkoly et al., 2021). is rise in research about dream and lucid dreaming ohers us novel and uniue perspectives to better understand what it is like to be conscious while sleeping. Lucid dreaming is the ability to become aware to become aware of dreaming while in a dream. References Casagrande, M., Violani, C., Lucidi, F., Buttinelli, E., \& Bertini, M. (1996). Variations in Sleep Mentation as a Function of Time of Night. International Journal of Neuroscience , 85 (1-2), 19-30. https://doi.org/10.3109/00207459608986348 ALIUS Bulletin n6 (2022) aliusresearch.org/bulletin Matthieu Koroma - What is it like to be asleep ? 4 Dement, W., \& Wolpert, E. A. (1958). e relation of eye movements, body motility, and external stimuli to dream content. Journal of Experimenal Psycholo , 55 (6), 543-553. https://doi.org/10.1037/h0040031 Dennett, D. C. (1976). Are dreams experiences?. e Philosophical Review , 85 (2), 151-171. https://doi.org/10.2307/2183728 Fazekas, P., Nemeth, G., \& Overgaard, M. (2019). White dreams are made of colours: What studying contentless dreams can teach about the neural basis of dreaming and conscious experiences. Sleep Medicine Reviews , 43 , 84-91. https://doi.org/10.1016/j.smrv.2018.10.005 Konkoly, K. R., Appel, K., Chabani, E., Mangiaruga, A., Gott, J., Mallett, R., Caughran, B., Witkowski, S., Whitmore, N. W., Mazurek, C. Y., Berent, J. B., Weber, F. D., Turker, B., Leu-Semenescu, S., Maranci, J.-B., Pipa, G., Arnulf, I., Oudiette, D., Dresler, M., \& Paller, K. A. (2021). Real-time dialogue between experimenters and dreamers during REM sleep. Current Biolo , S0960982221000592. https://doi.org/10.1016/j.cub.2021.01.026 LaBerge, S., \& Rheingold, H. (1990). Exploring the World of Lucid Dreaming. New York: Ballantine Books. Malcolm, N. (2017). Dreaming . Routledge. Oudiette, D., De Cock, V. C., Lavault, S., Leu, S., Vidailhet, M., \& Arnulf, I. (2009). Nonviolent elaborate behaviors may also occur in REM sleep behavior disorder. Neurolo , 72 (6), 551-557. https://doi.org/10.1212/01.wnl.0000341936.78678.3a Revonsuo, A. (2005). e self in dreams. In e lost self: Pathologies of the brain and identity (pp. 206-219). Oxford University Press Oxford. Siclari, F., LaRocue, J. J., Postle, B. R., \& Tononi, G. (2013). Assessing sleep consciousness within subjects using a serial awakening paradigm. Frontiers in Psycholo , 4 . https://doi.org/10.3389/fpsyg.2013.00542 Stickgold, R. (2000). Replaying the Game: Hypnagogic Images in Normals and Amnesics. Science , 290 (5490), 350-353. https://doi.org/10.1126/science.290.5490.350 Windt, J. M. (2013). Reporting dream experience: Why (not) to be skeptical about dream reports. Frontiers in Human Neuroscience , 7 . https://doi.org/10.3389/fnhum.2013.00708 ALIUS Bulletin n6 (2022) aliusresearch.org/bulletin Matthieu Koroma - What is it like to be asleep ? 5 Windt, J. M., Nielsen, T., \& ompson, E. (2016). Does Consciousness Disappear in Dreamless Sleep? Trends in Cognitive Sciences , 20 (12), 871-882. https://doi.org/10.1016/j.tics.2016.09.006 Ward, A. F., \& Wegner, D. M. (2013). Mind-blanking: When the mind goes away. Frontiers in psycholo , 4 , 650. https://doi.org/10.3389\%2Ffpsyg.2013.00650 Oudiette, D., Dodet, P., Ledard, N., Artru, E., Rachidi, I., Similowski, T., \& Arnulf, I. (2018). REM sleep respiratory behaviours match mental content in narcoleptic lucid dreamers. Scientific reports , 8 (1), 1-10. https://doi.org/10.1038/s41598-018-21067-9 ALIUS Bulletin n6 (2022) aliusresearch.org/bulletin\par
  \endgroup
}%
\end{document}
